\documentclass[11pt]{article}
\usepackage[utf8]{inputenc}
\usepackage[T2A]{fontenc}
\usepackage[russian]{babel}
% символы Unicode исходного текста Lean, чтобы цитировать его дословно
\DeclareUnicodeCharacter{2286}{\ensuremath{\subseteq}}
\DeclareUnicodeCharacter{2192}{\ensuremath{\to}}
\DeclareUnicodeCharacter{211D}{\ensuremath{\mathbb{R}}}
\DeclareUnicodeCharacter{2200}{\ensuremath{\forall}}
\DeclareUnicodeCharacter{2203}{\ensuremath{\exists}}
\DeclareUnicodeCharacter{2264}{\ensuremath{\le}}
\DeclareUnicodeCharacter{2227}{\ensuremath{\wedge}}
\DeclareUnicodeCharacter{2243}{\ensuremath{\simeq}}
\DeclareUnicodeCharacter{1D62}{\ensuremath{_i}}
\DeclareUnicodeCharacter{2194}{\ensuremath{\leftrightarrow}}
\usepackage{amsmath,amssymb,amsthm,mathtools}
\usepackage{tikz}
\usepackage[margin=1in]{geometry}
\setlength{\emergencystretch}{3em}
\usepackage[colorlinks=true,linkcolor=blue,citecolor=blue,urlcolor=blue]{hyperref}

\newtheorem{theorem}{Теорема}[section]
\newtheorem{lemma}[theorem]{Лемма}
\newtheorem{proposition}[theorem]{Предложение}
\newtheorem{corollary}[theorem]{Следствие}
\theoremstyle{definition}
\newtheorem{definition}[theorem]{Определение}
\theoremstyle{remark}
\newtheorem{remark}[theorem]{Замечание}
\newtheorem{example}[theorem]{Пример}

\newcommand{\R}{\mathbb{R}}
\newcommand{\Ort}{\mathrm{O}(3)}
\newcommand{\Fit}{\operatorname{Fit}}
\newcommand{\NC}{\operatorname{NC}}
\newcommand{\sgn}{\operatorname{sgn}}
\DeclarePairedDelimiter{\abs}{\lvert}{\rvert}

\title{Когда коробка помещается в коробку?\\ Явный критерий с формальным доказательством}
\author{Михаил Левит}
\date{27 сентября 2026~г.}

\begin{document}


\maketitle

\begin{abstract}
Мы даём необходимое и достаточное условие в явном виде для того, чтобы прямоугольный параллелепипед
с рёбрами $q_1,q_2,q_3\ge 0$ после произвольного движения поместился в прямоугольный параллелепипед
с рёбрами $p_1,p_2,p_3>0$. Условие представляет собой дизъюнкцию девяти явных проверок, каждая из
которых состоит из двух плоских тестов, выраженных формулой с одним квадратным корнем. Ключевой
геометрический факт: если коробка вообще помещается, то она помещается и в положении, в котором одно
из её рёбер перпендикулярно одной из осей ящика. Мы доказываем его, показывая, что в любом другом
положении все три ширины коробки можно одновременно уменьшить. Задача поставлена Гарнеттом в 1923
году; Джеррард и Вецель в 2004 году называли её частным случаем общей нерешённой задачи о вложении коробок в высших размерностях. Главная теорема проверена в системе
интерактивного доказательства Lean~4 с библиотекой Mathlib, вместе с аналогичным критерием строгого
вложения (без касания стенок). Он получается заменой всех нестрогих неравенств строгими, включая
неравенство, по которому выбирается ветвь формулы; замена только в неравенствах проверок даёт
неверные ответы.
Критерий конструктивен: если коробка помещается, он даёт явный поворот. Реализация критерия
отвечает на вопрос и строит положение коробки примерно за $20$ микросекунд, на три-четыре порядка
быстрее локального численного поиска по поворотам из многих начальных точек и без его ложных
ответов <<не помещается>> вблизи границы.
\end{abstract}

\section{Введение}

Пусть $P=[0,p_1]\times[0,p_2]\times[0,p_3]$, где $p_i>0$, и $Q=[0,q_1]\times[0,q_2]\times[0,q_3]$,
где $q_j\ge 0$. Будем говорить, что $Q$ \emph{помещается в} $P$, если $f(Q)\subseteq P$ для
некоторой изометрии $f$ пространства $\R^3$. Вопрос о том, когда это происходит, поставил
Ф.~М.~Гарнетт в 1923 году \cite{Garnett1923}; У.~Б.~Карвер \cite{Carver1925} дал частичный ответ,
приняв без доказательства, что в оптимальном положении все вершины $Q$ лежат на гранях $P$. Плоский
аналог, прямоугольник в прямоугольнике, решил Карвер в 1957 году \cite{Carver1957} (см. также Вецель
\cite{Wetzel2000}). Джеррард и Вецель \cite[с.~22]{JerrardWetzel2004} описывают трёхмерный вопрос как
частный случай <<общей нерешённой задачи о вложении коробок в высших размерностях>>. Из близких
результатов отметим диск в коробке \cite{AlexanderWetzel2005}, круговой цилиндр вдоль пространственной
диагонали коробки \cite{JSSW2006}, где критерий <<в духе Поста>> \cite{Post1993} назван желательным и
нерешённым, и алгоритмы вложения Мартина и Стивенсона \cite{MartinStephenson1989}. Олимпиадная задача
\cite{Kvant1999} доказывает необходимое условие $q_1+q_2+q_3\le p_1+p_2+p_3$.

В настоящей работе такой критерий дан для коробок. Он является дизъюнкцией явных проверок, подобно
тому как критерий Поста для треугольника в треугольнике \cite{Post1993} является дизъюнкцией 18
неравенств.

\subsection*{Критерий}
Для $x\ge y\ge0$ и $s\ge 0$ положим
\begin{equation}\label{eq:H}
H(x,y;s)=\begin{cases}
y, & x\le s,\\[2pt]
\min\Bigl(x,\ \dfrac{2xys+(x^2-y^2)\sqrt{x^2+y^2-s^2}}{x^2+y^2}\Bigr), & x>s,
\end{cases}
\end{equation}
а для произвольных $u,v\ge0$ положим $H(u,v;s)=H(\max(u,v),\min(u,v);s)$. По
лемме~\ref{lem:strip} при $\min(u,v)\le s$ число $H(u,v;s)$ есть наименьшая возможная
протяжённость вдоль полосы ширины $s$ прямоугольника $u\times v$, поставленного в эту полосу. Для
$a,b,u,v\ge 0$ положим
\[
\Fit(a,b;u,v)\ :\Longleftrightarrow\ \min(u,v)\le a\ \text{ и }\ H(u,v;a)\le b .
\]
Это условие означает, что прямоугольник $u\times v$ помещается в прямоугольник $a\times b$
(замечание~\ref{rem:carver}).

\begin{definition}[Критерий $\NC$]\label{def:NC}
Пусть $(i,i',i'')$~--- перестановка осей $(1,2,3)$ ящика $P$, а $(j,k,l)$~--- перестановка рёбер
$(1,2,3)$ коробки $Q$. \emph{Проверка} $\NC_{(i,i',i''),(j,k,l)}$ состоит в том, что
\[
\min(q_k,q_l)\le p_i\quad\text{и}\quad \Fit\bigl(p_{i'},p_{i''};\,q_j,\,h^*\bigr),\qquad
h^*:=H(q_k,q_l;p_i).
\]
Условие $\NC(p,q)$ выполнено, если выполнена хотя бы одна проверка. Поскольку $H$ симметрична по
первым двум аргументам, имеется $9$ вариантов выбора $(i,j)$ и для каждого два порядка осей <<пола>>
$(i',i'')$.
\end{definition}

Геометрически проверка описывает положения, в которых ребро $j$ коробки $Q$ горизонтально, то есть
перпендикулярно оси $i$ ящика $P$ (<<высоте>>). Грань $q_k\times q_l$ при этом стоит в вертикальной
полосе высоты $p_i$; её ортогональная проекция на пол~--- отрезок длины не меньше $h^*$, и прямоугольник
$q_j\times h^*$ должен поместиться в пол $p_{i'}\times p_{i''}$ (рис.~\ref{fig:box}).

\begin{figure}[t]
\centering
\begin{tikzpicture}[line join=round]
\fill[black!6] (0.000,0.000) -- (3.539,1.163) -- (1.115,2.788) -- (-2.424,1.625) -- cycle;
\fill[black!18] (-0.644,0.432) -- (-0.327,0.401) -- (0.866,0.285) -- (1.778,2.344) -- (0.267,2.491) -- cycle;
\draw[black!60] (0.000,0.000) -- (0.000,1.695);
\draw[black!60] (0.000,0.000) -- (-2.424,1.625);
\draw[black!60] (0.000,0.000) -- (3.539,1.163);
\draw[black!60] (0.000,1.695) -- (-2.424,3.320);
\draw[black!60] (0.000,1.695) -- (3.539,2.859);
\draw[black!60] (-2.424,1.625) -- (-2.424,3.320);
\draw[densely dashed, black!45] (-2.424,1.625) -- (1.115,2.788);
\draw[black!60] (-2.424,3.320) -- (1.115,4.484);
\draw[black!60] (3.539,1.163) -- (3.539,2.859);
\draw[densely dashed, black!45] (3.539,1.163) -- (1.115,2.788);
\draw[black!60] (3.539,2.859) -- (1.115,4.484);
\draw[densely dashed, black!45] (1.115,2.788) -- (1.115,4.484);
\filldraw[fill=blue!22, draw=blue!50!black, fill opacity=0.9] (1.778,2.551) -- (0.585,4.156) -- (0.267,3.980) -- (1.461,2.375) -- cycle;
\filldraw[fill=blue!35, draw=blue!50!black, fill opacity=0.9] (0.866,0.492) -- (1.778,2.551) -- (1.461,2.375) -- (0.549,0.316) -- cycle;
\filldraw[fill=blue!12, draw=blue!50!black, fill opacity=0.9] (0.549,0.316) -- (1.461,2.375) -- (0.267,3.980) -- (-0.644,1.920) -- cycle;
\filldraw[fill=blue!12, draw=blue!50!black, fill opacity=0.9] (0.866,0.492) -- (1.778,2.551) -- (0.585,4.156) -- (-0.327,2.096) -- cycle;
\filldraw[fill=blue!35, draw=blue!50!black, fill opacity=0.9] (-0.327,2.096) -- (0.585,4.156) -- (0.267,3.980) -- (-0.644,1.920) -- cycle;
\filldraw[fill=blue!22, draw=blue!50!black, fill opacity=0.9] (0.866,0.492) -- (-0.327,2.096) -- (-0.644,1.920) -- (0.549,0.316) -- cycle;
\draw[black!60] (0.000,0.000) -- (0.000,1.695);
\draw[black!60] (0.000,1.695) -- (-2.424,3.320);
\draw[black!60] (0.000,1.695) -- (3.539,2.859);
\end{tikzpicture}
\caption{Коробка $2{,}8\times1{,}3\times0{,}25$ в ящике высоты $1{,}2$ с полом $2{,}7\times2{,}64$ в положении, которое описывает одна из проверок критерия: длинное ребро горизонтально, грань $1{,}3\times0{,}25$ наклонена в вертикальной полосе, а ортогональная проекция коробки на пол (тень)~--- прямоугольник $2{,}8\times h^*$, повёрнутый на полу. Поворот вычислен как в доказательстве предложения~\ref{prop:suff}.}\label{fig:box}
\end{figure}

\begin{theorem}[Главная теорема]\label{thm:main}
Пусть $p_1,p_2,p_3>0$ и $q_1,q_2,q_3\ge 0$. Коробка $Q$ помещается в ящик $P$ тогда и только тогда,
когда выполнено $\NC(p,q)$.
\end{theorem}

Достаточность конструктивна (раздел~\ref{sec:suff}): каждая выполненная проверка даёт явный поворот.
Необходимость опирается на следующую теорему (раздел~\ref{sec:descent}), которая является главным
новым элементом. Здесь $w_i(R)=\sum_j q_j\abs{R_{ij}}$~--- ширина повёрнутой коробки вдоль оси $i$
(раздел~\ref{sec:widths}).

\begin{theorem}[Спуск]\label{thm:Z}
Пусть $q_1,q_2,q_3>0$ и матрица $R\in\Ort$ не имеет нулевых элементов. Тогда существует $R'\in\Ort$,
для которой $w_i(R')<w_i(R)$ при $i=1,2,3$.
\end{theorem}

\subsection*{Формальная проверка}
Теорема~\ref{thm:main} формализована и проверена в Lean~4 с Mathlib; см. раздел~\ref{sec:lean}.
Формальная формулировка буквально использует приведённые выше формулы; её смысл зависит только от
определений коробки и вложения.

\section{Сведение к ширинам}\label{sec:widths}

Всякая изометрия $\R^3$ имеет вид $x\mapsto Rx+t$, где $R\in\Ort$ (Мазур--Улам). Столбец $j$
матрицы $R$~--- направление ребра $j$ перемещённой коробки, а $R_{ij}$~--- косинус угла между этим
ребром и осью $i$ ящика $P$.

\begin{lemma}\label{lem:widths}
Пусть $R$~--- произвольная вещественная матрица $3\times3$. Некоторый сдвиг множества $RQ$ лежит в $P$
тогда и только тогда, когда
\[
w_i(R):=\sum_{j=1}^3 q_j\abs{R_{ij}}\le p_i\qquad(i=1,2,3).
\]
Следовательно, $Q$ помещается в $P$ тогда и только тогда, когда $w(R)\le p$ для некоторой $R\in\Ort$.
\end{lemma}

\begin{proof}
При $x\in Q$ координата $i$ точки $Rx=\sum_j R_{ij}x_j$ пробегает отрезок длины
$\sum_j q_j\abs{R_{ij}}$, концы которого достигаются в вершинах $Q$. Сдвиги вдоль разных осей
независимы.
\end{proof}

То же рассуждение на плоскости показывает, что прямоугольник $u\times v$ помещается в прямоугольник
$a\times b$ тогда и только тогда, когда $uc+vd\le a$ и $ud+vc\le b$ для некоторых $c,d\ge0$ с
$c^2+d^2=1$.

\section{Прямоугольник в полосе}\label{sec:strip}

Пусть $x\ge y\ge 0$ и $s\ge0$. Для $c,d\ge0$ с $c^2+d^2=1$ положим $A=xc+yd$ (протяжённость
повёрнутого прямоугольника $x\times y$ поперёк полосы) и $B=xd+yc$ (его протяжённость вдоль полосы).

\begin{lemma}[Лемма о полосе]\label{lem:strip}
\begin{enumerate}
\item[(a)] Если $y\le s$, то существуют $c,d\ge 0$ с $c^2+d^2=1$, $A\le s$ и $B\le H(x,y;s)$.
\item[(b)] Если $c,d\ge0$, $c^2+d^2=1$ и $A\le s$, то $y\le s$ и $H(x,y;s)\le B$.
\end{enumerate}
То же верно для произвольных $u,v\ge0$ вместо $x,y$ с $H(u,v;s)$, $A=uc+vd$, $B=ud+vc$ (при $u<v$
следует поменять местами $c$ и $d$).
\end{lemma}

\begin{proof}
При $x=0$ всё тривиально, поэтому пусть $r=\sqrt{x^2+y^2}>0$ и $\alpha\in[0,\pi/4]$ задан условиями
$\cos\alpha=x/r$, $\sin\alpha=y/r$. Запишем $(c,d)=(\cos\theta,\sin\theta)$, $\theta\in[0,\pi/2]$.
Тогда
\[
A(\theta)=r\cos(\theta-\alpha),\qquad B(\theta)=r\sin(\theta+\alpha).
\]
При $\theta\in[0,\pi/2]$ имеем $\theta-\alpha\in[-\pi/2,\pi/2]$ и $\theta+\alpha\in[0,\pi]$, поэтому
$A$ и $B$ вогнуты на $[0,\pi/2]$; значит, на любом отрезке они достигают минимума на конце. Кроме того,
$A(0)=B(\pi/2)=x$ и $A(\pi/2)=B(0)=y$.

По вогнутости $A\ge\min(x,y)=y$, что доказывает первое утверждение (b). Если $x\le s$, то угол
$\theta=0$ допустим, $B(0)=y$, и всюду $B\ge\min(B(0),B(\pi/2))=y$; это доказывает (a) и (b).

Пусть $y\le s<x$. На $[0,\alpha]$ функция $A$ возрастает от $x$ до $r$, так что там $A>s$; на
$[\alpha,\pi/2]$ она убывает от $r$ до $y\le s$. Поэтому допустимые углы образуют отрезок
$[\theta^*,\pi/2]$, где $A(\theta^*)=s$, и минимум вогнутой функции $B$ на нём равен
$\min(B(\theta^*),B(\pi/2))=\min(B(\theta^*),x)$. Остаётся вычислить $B(\theta^*)$. При
$\varphi=\theta^*-\alpha\in[0,\pi/2]$ имеем $\cos\varphi=s/r$, $\sin\varphi=\sqrt{r^2-s^2}/r$ и
\[
\cos\theta^*=\frac{xs-y\sqrt{r^2-s^2}}{r^2},\qquad
\sin\theta^*=\frac{ys+x\sqrt{r^2-s^2}}{r^2},\qquad
B(\theta^*)=\frac{2xys+(x^2-y^2)\sqrt{r^2-s^2}}{r^2}.
\]
Итак, $\min\{B:A\le s\}=H(x,y;s)$, откуда следуют (a) (берём конец, на котором достигается минимум)
и (b).
\end{proof}

\begin{remark}\label{rem:carver}
По плоскому варианту леммы~\ref{lem:widths} и лемме~\ref{lem:strip} условие $\Fit(a,b;u,v)$
выполнено тогда и только тогда, когда прямоугольник $u\times v$ помещается в прямоугольник
$a\times b$. Карвер \cite{Carver1957} дал другую явную форму того же условия: при $u\ge v$, $a\ge b$,
$u>a$, $v\le b$ оно имеет вид
$\bigl(\tfrac{a+b}{u+v}\bigr)^2+\bigl(\tfrac{a-b}{u-v}\bigr)^2\ge 2$. Форма Карвера нами не
используется. (Численно обе формы совпадают на $2\cdot10^5$ случайных прямоугольниках; см.
\texttt{tools/box\_fit.py}.)
\end{remark}

\section{Достаточность}\label{sec:suff}

\begin{proposition}\label{prop:suff}
Если выполнена какая-либо проверка $\NC(p,q)$, то $Q$ помещается в $P$, причём поворот можно
выписать явно.
\end{proposition}

\begin{proof}
Перенумеровав оси и рёбра (что переставляет строки и столбцы $R$), можно считать, что выполнена
проверка $\NC_{(1,2,3),(1,2,3)}$: $\min(q_2,q_3)\le p_1$ и $\Fit(p_2,p_3;q_1,h^*)$, где
$h^*=H(q_2,q_3;p_1)$. По лемме~\ref{lem:strip}(a) существуют $c_1,d_1\ge0$, $c_1^2+d_1^2=1$, с
\[
q_2c_1+q_3d_1\le p_1,\qquad q_2d_1+q_3c_1\le h^*,
\]
в частности $h^*\ge0$. Применяя лемму~\ref{lem:strip}(a) к прямоугольнику $q_1\times h^*$ в полосе
ширины $p_2$, получаем $c_2,d_2\ge 0$, $c_2^2+d_2^2=1$, с
\[
q_1c_2+h^*d_2\le p_2,\qquad q_1d_2+h^*c_2\le H(q_1,h^*;p_2)\le p_3 .
\]
Положим
\[
R=\begin{pmatrix}0&c_1&-d_1\\ c_2&-d_1d_2&-c_1d_2\\ d_2&d_1c_2&c_1c_2\end{pmatrix}.
\]
Её столбцы $f_1=(0,c_2,d_2)$, $f_2=c_1e_1+d_1n$, $f_3=-d_1e_1+c_1n$, где $n=(0,-d_2,c_2)$,
ортонормированы, так что $R\in\Ort$. Её ширины равны
\begin{align*}
w_1&=q_2c_1+q_3d_1\le p_1,\\
w_2&=q_1c_2+d_2(q_2d_1+q_3c_1)\le q_1c_2+d_2h^*\le p_2,\\
w_3&=q_1d_2+c_2(q_2d_1+q_3c_1)\le q_1d_2+c_2h^*\le p_3,
\end{align*}
и остаётся применить лемму~\ref{lem:widths}.
\end{proof}

\section{Положения с горизонтальным ребром}\label{sec:zero}

\begin{proposition}\label{prop:zero}
Пусть $q\ge0$, $p>0$, и матрица $R\in\Ort$ удовлетворяет условиям $w(R)\le p$ и $R_{ij}=0$ для
некоторых $i,j$. Тогда выполнено $\NC(p,q)$.
\end{proposition}

\begin{proof}
Перенумеровав, считаем $R_{11}=0$. Для ортогональной матрицы $R=(\det R)\operatorname{cof}(R)$ и
$\det R=\pm1$, поэтому каждое $\abs{R_{ab}}$ равно модулю соответствующего алгебраического
дополнения. Так как $R_{11}=0$, получаем
\[
\abs{R_{22}}=\abs{R_{13}}\abs{R_{31}},\quad \abs{R_{23}}=\abs{R_{12}}\abs{R_{31}},\quad
\abs{R_{32}}=\abs{R_{13}}\abs{R_{21}},\quad \abs{R_{33}}=\abs{R_{12}}\abs{R_{21}}.
\]
Положим $c=\abs{R_{12}}$, $d=\abs{R_{13}}$, $c'=\abs{R_{21}}$, $d'=\abs{R_{31}}$; тогда
$c^2+d^2=1$ (первая строка) и $c'^2+d'^2=1$ (первый столбец). Обозначив $B:=q_2d+q_3c$, получаем
\[
w_1=q_2c+q_3d,\qquad w_2=q_1c'+d'B,\qquad w_3=q_1d'+c'B .
\]
Из $w_1\le p_1$ и леммы~\ref{lem:strip}(b) следует $\min(q_2,q_3)\le p_1$ и
$h^*:=H(q_2,q_3;p_1)\le B$; кроме того, $h^*\ge0$ по лемме~\ref{lem:strip}(a). Значит,
$q_1c'+d'h^*\le w_2\le p_2$ и $q_1d'+c'h^*\le w_3\le p_3$. Лемма~\ref{lem:strip}(b), применённая к
прямоугольнику $q_1\times h^*$ в полосе ширины $p_2$ с направлением $(c',d')$, даёт
$\min(q_1,h^*)\le p_2$ и $H(q_1,h^*;p_2)\le q_1d'+h^*c'\le p_3$, то есть $\Fit(p_2,p_3;q_1,h^*)$.
Итак, выполнена проверка $\NC_{(1,2,3),(1,2,3)}$.
\end{proof}

\section{Спуск: доказательство теоремы~\ref{thm:Z}}\label{sec:descent}

Всюду в этом разделе $q_1,q_2,q_3>0$ и $R\in\Ort$ не имеет нулевых элементов. Через $e_1,e_2,e_3$
обозначен стандартный базис, через $u\cdot v$~--- скалярное произведение, через $u\times v$~---
векторное.

\subsection{Линеаризация}
Пусть $s_{ij}=\sgn R_{ij}\in\{\pm1\}$, $g_i=(s_{i1}q_1,s_{i2}q_2,s_{i3}q_3)$ и $a_i=Rg_i\in\R^3$.
Обозначим $A_{ik}=(a_i)_k$. Тогда $A_{ii}=\sum_j R_{ij}s_{ij}q_j=w_i(R)$. Если $C\in\Ort$ и знаки
элементов $CR$ совпадают со знаками элементов $R$, то
\begin{equation}\label{eq:lin}
w_i(CR)=\sum_j q_j s_{ij}(CR)_{ij}=(Ca_i)_i .
\end{equation}

\begin{lemma}[B1]\label{lem:B1}
$\abs{A_{ik}}<A_{kk}$ при $i\neq k$.
\end{lemma}
\begin{proof}
$\abs{A_{ik}}=\abs{\sum_j s_{ij}q_jR_{kj}}\le\sum_jq_j\abs{R_{kj}}=A_{kk}$. Равенство означало бы,
что все $s_{ij}R_{kj}$ ($q_j>0$, $R_{kj}\ne0$) одного знака, а тогда все произведения
$R_{ij}R_{kj}$ ненулевые и одного знака, что противоречит ортогональности строк $i$ и $k$.
\end{proof}

\subsection{Поворот Кэли}
Для $v\in\R^3$ положим $[v]_\times x=v\times x$ и
\[
C(v)=\frac{(1-\abs{v}^2)I+2[v]_\times+2vv^{\mathsf T}}{1+\abs{v}^2}.
\]
Это поворот, задаваемый кватернионом $(1,v)$; непосредственно проверяется, что
$C(v)C(v)^{\mathsf T}=I$. Кроме того, $C(0)=I$ и $C$ непрерывна. Для $a\in\R^3$ положим
\[
b_i(a)=a\times e_i,\qquad P_i(a;v)=v_i\,(v\cdot a)-\abs{v}^2a_i,\qquad D_i(a;v)=b_i(a)\cdot v+P_i(a;v).
\]
Так как $e_i\cdot(v\times a)=v\cdot(a\times e_i)$, прямое вычисление даёт
\begin{equation}\label{eq:cay}
(C(v)a)_i-a_i=\frac{2}{1+\abs{v}^2}\,D_i(a;v).
\end{equation}
Заметим, что $P_i(a;\cdot)$~--- квадратичная форма: $P_i(a;tv)=t^2P_i(a;v)$. Именно поэтому мы
пользуемся параметризацией Кэли: по \eqref{eq:cay} изменение каждой ширины с точностью до
положительного множителя есть явный многочлен второй степени от $v$, так что доступны и члены
первого, и члены второго порядка.

\subsection{Спуск вдоль кривой}
Будем поворачивать вдоль $v(t)=t(h+tz)$ при малых $t>0$. Тогда
\begin{equation}\label{eq:curve}
D_i(a;v(t))=t\,b_i(a)\cdot h+t^2\bigl(b_i(a)\cdot z+P_i(a;h+tz)\bigr).
\end{equation}

\begin{lemma}\label{lem:order}
Пусть $a,h,z\in\R^3$ и $b=b_i(a)$.
\begin{enumerate}
\item[(a)] Если $b\cdot h<0$, либо $b\cdot h=0$ и $b\cdot z+P_i(a;h)<0$, то $D_i(a;v(t))<0$ при
всех достаточно малых $t>0$.
\item[(b)] Если $b=0$, $a_i>0$ и все компоненты $z$ ненулевые, то $D_i(a;v(t))<0$ при всех
достаточно малых $t>0$.
\end{enumerate}
\end{lemma}
\begin{proof}
(a) По \eqref{eq:curve} $D_i/t=b\cdot h+t(b\cdot z+P_i(a;h+tz))\to b\cdot h$ при $t\to0^+$; во втором
случае $D_i/t^2=b\cdot z+P_i(a;h+tz)\to b\cdot z+P_i(a;h)$.
(b) Условие $b=a\times e_i=0$ означает $a=a_ie_i$; тогда
$P_i(a;u)=u_i^2a_i-\abs{u}^2a_i=-a_i\sum_{m\ne i}u_m^2$. При малых $t>0$ каждая компонента
$h_m+tz_m$ вектора $u=h+tz$ ненулевая (при $h_m\ne0$ по непрерывности, при $h_m=0$ потому, что
$z_m\ne0$). Следовательно, $D_i(a;v(t))=t^2P_i(a;h+tz)<0$.
\end{proof}

\subsection{Ключевое неравенство}

\begin{lemma}\label{lem:key}
Пусть $a_1,a_2,a_3\in\R^3$ удовлетворяют условию $\abs{A_{ik}}<A_{kk}$ при $i\ne k$, где
$A_{ik}=(a_i)_k$. Пусть $\nu\in\R^3$, $\nu\ge0$, имеет не менее двух положительных компонент и
$\sum_i\nu_i\,b_i(a_i)=0$. Тогда $\sum_i\nu_iP_i(a_i;h)<0$ для любого $h\ne0$.
\end{lemma}

\begin{proof}
Так как $a\times e_1=(0,a_3,-a_2)$, $a\times e_2=(-a_3,0,a_1)$, $a\times e_3=(a_2,-a_1,0)$, условие
$\sum_i\nu_ib_i(a_i)=0$ в точности означает, что
\[
\nu_iA_{ik}=\nu_kA_{ki}\qquad(i\ne k).
\]
Положим $d_i=\nu_iA_{ii}\ge0$ и $m_{ik}=\nu_iA_{ik}=\nu_kA_{ki}$. Раскрывая скобки и пользуясь
симметрией, получаем
\[
\sum_i\nu_iP_i(a_i;h)=\sum_i\nu_i\Bigl(h_i\sum_kh_kA_{ik}-\abs{h}^2A_{ii}\Bigr)=-\sum_{i<k}B_{ik},
\qquad B_{ik}=d_ih_k^2+d_kh_i^2-2m_{ik}h_ih_k .
\]
По условию
$m_{ik}^2=\nu_i\nu_kA_{ik}A_{ki}\le\nu_i\nu_k\abs{A_{ik}}\abs{A_{ki}}\le\nu_i\nu_kA_{kk}A_{ii}=d_id_k$,
причём неравенство строгое при $\nu_i\nu_k>0$.

Каждое $B_{ik}\ge0$: если $d_i=0$, то $m_{ik}=0$ и $B_{ik}=d_kh_i^2$; если $d_i>0$, то
$d_iB_{ik}=(d_ih_k-m_{ik}h_i)^2+(d_id_k-m_{ik}^2)h_i^2\ge0$.

Некоторое $B_{ik}$ положительно. Выберем $c$ с $h_c\ne0$ и, пользуясь тем, что у $\nu$ две
положительные компоненты, индекс $i\ne c$ с $\nu_i>0$, так что $d_i>0$. Если $d_c=0$, то
$m_{ic}=0$ и $B_{ic}=d_ih_c^2>0$. Если $d_c>0$, то $\nu_c>0$, поэтому $m_{ic}^2<d_id_c$, и
$d_iB_{ic}=(d_ih_c-m_{ic}h_i)^2+(d_id_c-m_{ic}^2)h_i^2$ положительно как при $h_i\ne0$, так и
при $h_i=0$.
\end{proof}

\subsection{Выбор кривой}
Идея такова. Либо некоторое направление $h$ уменьшает все ширины в первом порядке, либо члены первого
порядка $b_i\cdot h$ заблокированы неотрицательным соотношением $\sum_i\nu_ib_i=0$; во втором случае
лемма~\ref{lem:key} показывает, что ширины уменьшают члены второго порядка.
Мы используем теорему Моцкина об альтернативах \cite{Motzkin1936,Schrijver1986} в следующей форме:
для векторов $b_i\in\R^3$ и чисел $c_i$, $i\in K$, либо система $b_i\cdot z<c_i$ ($i\in K$)
совместна, либо существует $\mu\ge0$, $\mu\ne0$, сосредоточенный на $K$, с $\sum_i\mu_ib_i=0$ и
$\sum_i\mu_ic_i\le0$.

\begin{lemma}\label{lem:choice}
Пусть $a_i$ те же, что в линеаризации, $b_i=b_i(a_i)$ и $K=\{i:b_i\ne0\}$. Существуют
$h,z\in\R^3$, такие что все компоненты $z$ ненулевые и для каждого $i\in K$ либо $b_i\cdot h<0$,
либо $b_i\cdot h=0$ и $b_i\cdot z+P_i(a_i;h)<0$.
\end{lemma}

\begin{proof}
Применим теорему Моцкина с $c=0$. Если система $b_i\cdot h<0$, $i\in K$, имеет решение $h$, возьмём
$z=(1,1,1)$. Иначе существует $\mu\ne0$, $\mu\ge0$, с $\sum\mu_ib_i=0$ (напомним, что $b_i=0$ при
$i\notin K$), поэтому матрица со строками $b_1,b_2,b_3$ вырождена и найдётся $h\ne0$ с
$b_i\cdot h=0$ для всех $i$. Применим теорему Моцкина ещё раз, с $c_i=-P_i(a_i;h)$.

Если существует $z_0$ с $b_i\cdot z_0+P_i(a_i;h)<0$ при $i\in K$, то $z=z_0+\delta(1,1,1)$ при малом
$\delta>0$ сохраняет эти строгие неравенства и имеет ненулевые компоненты.

Иначе существует $\nu\ge0$, $\nu\ne0$, сосредоточенный на $K$, с $\sum\nu_ib_i=0$ и
$\sum\nu_iP_i(a_i;h)\ge0$. Если бы у $\nu$ была лишь одна положительная компонента $\nu_i$, то из
$\nu_ib_i=0$ следовало бы $b_i=0$ вопреки $i\in K$. Значит, у $\nu$ две положительные компоненты, и
лемма~\ref{lem:key} (её условие выполнено по лемме~\ref{lem:B1}) даёт $\sum\nu_iP_i(a_i;h)<0$~---
противоречие.
\end{proof}

\subsection{Доказательство теоремы~\ref{thm:Z}}
Возьмём $h,z$ из леммы~\ref{lem:choice} и положим $R(t)=C(v(t))R$. При всех достаточно малых $t>0$:
\begin{itemize}
\item $R(t)\in\Ort$ и знаки её элементов совпадают со знаками элементов $R$, так как $R(t)\to R$, а
у $R$ нет нулевых элементов;
\item $D_i(a_i;v(t))<0$ для каждого $i$: при $i\in K$ по лемме~\ref{lem:order}(a), при $i\notin K$~---
по лемме~\ref{lem:order}(b), поскольку $(a_i)_i=w_i(R)>0$.
\end{itemize}
По \eqref{eq:lin} и \eqref{eq:cay}
$w_i(R(t))-w_i(R)=(C(v(t))a_i)_i-(a_i)_i=\frac{2}{1+\abs{v(t)}^2}D_i(a_i;v(t))<0$ при $i=1,2,3$.
\qed

\section{Доказательство главной теоремы}\label{sec:proof}

Достаточность~--- это предложение~\ref{prop:suff}. Для необходимости, в силу
предложения~\ref{prop:zero}, достаточно показать: \emph{если $Q$ помещается в $P$, то $w(R)\le p$
для некоторой $R\in\Ort$ с нулевым элементом.}

\emph{Случай $q_1,q_2,q_3>0$.} Функция $F(R)=\max_iw_i(R)/p_i$ непрерывна на компактной группе
$\Ort$ и достигает минимума в некоторой точке $R^*$; поскольку $Q$ помещается, $F(R^*)\le1$, то есть
$w(R^*)\le p$. Если бы у $R^*$ не было нулевых элементов, теорема~\ref{thm:Z} дала бы $R'$ с
$w_i(R')<w_i(R^*)$ для всех $i$, а значит, $F(R')<F(R^*)$. Следовательно, у $R^*$ есть нулевой
элемент.

\emph{Общий случай $q\ge0$.} Пусть $w(R_0)\le p$, и предположим, что ни одна $R$ из компактного
множества $Z=\{R\in\Ort:\ R_{ij}=0\text{ для некоторых }i,j\}$ не удовлетворяет условию $w(R)\le p$.
Тогда непрерывная функция $\varphi(R)=\max_i(w_i(R)-p_i)$ положительна на $Z$ и имеет на нём минимум
$m>0$. Пусть $\delta=m/6$, $q'=q+\delta(1,1,1)>0$ и $p'=p+3\delta(1,1,1)$. Так как
$\sum_j\abs{R_{ij}}\le3$ при $R\in\Ort$, имеем
$w_i(R_0;q')=w_i(R_0;q)+\delta\sum_j\abs{(R_0)_{ij}}\le p_i'$, то есть $q'$ помещается в $p'$. По
предыдущему случаю найдётся $R_1\in Z$ с $w(R_1;q')\le p'$, откуда
$w_i(R_1;q)\le w_i(R_1;q')\le p_i+3\delta$ и $\varphi(R_1)\le3\delta<m$~--- противоречие. \qed

\section{Примеры}\label{sec:examples}

\begin{example}[Стержень]
Пусть $q=(L,0,0)$. Возьмём проверку с высотой $i=1$, горизонтальным ребром $j=2$ и $(k,l)=(1,3)$.
Тогда $h^*=H(L,0;p_1)=\sqrt{\max(L^2-p_1^2,0)}$, и условие $\Fit(p_2,p_3;0,h^*)$ сводится к
$h^{*2}\le p_2^2+p_3^2$. Итак, проверка выполнена тогда и только тогда, когда
$L^2\le p_1^2+p_2^2+p_3^2$; обратно, более длинный стержень поместиться не может, поскольку
$\sqrt{p_1^2+p_2^2+p_3^2}$~--- диаметр $P$. Стержень помещается ровно по пространственной диагонали,
и в этом положении есть ребро (нулевой длины), перпендикулярное одной из осей.
\end{example}

\begin{example}[Квадрат в кубе]
Наибольший квадрат, помещающийся в единичный куб, имеет сторону $a=3\sqrt2/4$ (Нивланд; опубликовано
ван Свинденом \cite{vanSwinden1816}). Для $q=(a,a,0)$ и $p=(1,1,1)$ возьмём $j=1$, $(k,l)=(2,3)$:
$h^*=H(a,0;1)=\sqrt{a^2-1}=\sqrt2/4$. Для пола $x^2=9/8$, $y^2=1/8$, $x^2+y^2=5/4$,
$\sqrt{x^2+y^2-1}=1/2$, $2xy=3/4$, так что вторая ветвь \eqref{eq:H} равна
$(3/4+1\cdot\tfrac12)/(5/4)=1$ и $H(a,h^*;1)=\min(a,1)=1$: проверка выполнена с равенством, в
согласии с результатом Нивланда.
\end{example}

\begin{example}[Классических необходимых условий недостаточно]
Для $p=(1,1,1)$ и $q=(1{,}1;\,1{,}1;\,0{,}01)$ сумма рёбер, сумма квадратов рёбер, площадь
поверхности и объём у $Q$ меньше, однако $Q$ не помещается: её грань $1{,}1\times1{,}1$ дала бы в
кубе квадрат со стороной больше $3\sqrt2/4$.
\end{example}

\section{Вычисления}\label{sec:computation}

Критерий состоит не более чем из $36$ вычислений по формуле~\eqref{eq:H} и сравнений, так что он
проверяется за постоянное время. Если проверка выполнена, доказательство
предложения~\ref{prop:suff} даёт поворот явно: два угла из леммы~\ref{lem:strip}(a), один для грани в
вертикальной полосе, другой для прямоугольника на полу. Эталонная реализация
\texttt{tools/box\_fit.py} в репозитории (раздел~\ref{sec:lean}) возвращает матрицу поворота $R$ с
$\det R=1$ и сдвиг $t$, для которых $x\mapsto Rx+t$ переводит $Q$ в $P$; её самопроверка проверяет
около $2\cdot10^4$ случайных положений (ортогональность $R$ и то, что все вершины лежат в $P$).

Для сравнения мы использовали численный поиск, который минимизирует $\max_iw_i(R)/p_i$ по поворотам
методом Нелдера--Мида из случайных начальных точек и отвечает <<помещается>>, если из какой-то точки
достигнуто значение $\le1$. На $200$ случайных парах вблизи границы (коробка масштабирована с
множителем $1\pm e$ от наибольшего помещающегося размера, $10^{-4}\le e\le10^{-2}$) критерий тратит около
$24$ микросекунд на пару, построение положения~--- около $20$ микросекунд (Python, одно ядро
настольного компьютера), а численный поиск~--- $78$, $228$ и $517$ миллисекунд на пару при $5$, $20$ и
$50$ начальных точках, т.\,е. он медленнее в $3\cdot10^3$--$2\cdot10^4$ раз и даёт соответственно $27$,
$5$ и $0$ неверных ответов. Все его ошибки~--- ответ <<не помещается>> для коробок, которые помещаются:
локальный поиск может подтвердить, что коробка помещается, но то, что он не нашёл положения, ничего
не доказывает. Строгий отрицательный ответ требует глобального метода с гарантированными оценками,
например перебора по поворотам методом ветвей и границ, что значительно дороже; критерий даёт оба
ответа точно. Скрипт замера~--- \texttt{tools/benchmark.py}.

\section{Строгое вложение}\label{sec:strict}

\begin{remark}\label{rem:strict}
Будем говорить, что $Q$ \emph{помещается строго} в $P$, если $f(Q)$ лежит в открытом параллелепипеде
$(0,p_1)\times(0,p_2)\times(0,p_3)$ для некоторой изометрии $f$, то есть не касаясь границы $P$.
Пусть $H^\circ$ определена как $H$ в \eqref{eq:H}, но с разбором случаев $x<s$ / $x\ge s$ вместо
$x\le s$ / $x>s$, и пусть $\NC^\circ$~--- критерий из определения~\ref{def:NC}, в котором все
неравенства строгие, а $H$ заменена на $H^\circ$. Тогда при $p>0$ и $q\ge0$:
\begin{itemize}
\item $Q$ помещается строго в $P$ тогда и только тогда, когда $\NC(p-\varepsilon(1,1,1),q)$ выполнено
при некотором $0<\varepsilon<\min_ip_i$;
\item $Q$ помещается строго в $P$ тогда и только тогда, когда выполнено $\NC^\circ(p,q)$.
\end{itemize}
Полные формальные доказательства находятся в \texttt{lean/Boxes/Strict.lean} (теоремы
\texttt{fitsStrict\_iff\_shrink} и \texttt{fitsStrict\_iff\_NCs}); они повторяют
разделы~\ref{sec:widths}--\ref{sec:proof} со строгими неравенствами.

Изменение разбора случаев существенно, так как $H(x,y;s)$ имеет скачок при $s=x$: грань, длинная
сторона которой равна высоте, может стоять вертикально, касаясь потолка, но если касаться нельзя, её
не удаётся слегка наклонить (малый наклон увеличивает её высоту). Например, пусть $p=(1,1,1)$ и
$q=(1;\,0{,}5;\,0{,}5)$. Если сделать строгими только неравенства проверок, оставив $H$ прежней, одна из проверок
выполнена: при
горизонтальном ребре $0{,}5$ имеем $h^*=H(1;0{,}5;1)=0{,}5$, и $0{,}5\times0{,}5$ строго помещается в
$1\times1$. Но $\NC^\circ$ не выполнен: $H^\circ(1;0{,}5;1)=\min\bigl(1;\,0{,}5\cdot2{,}75/1{,}25\bigr)=1$,
а каждая проверка требует $H^\circ(1;0{,}5;1)<1$ (для пола, если горизонтально ребро $1$, и для $h$, а
затем для пола, если горизонтально ребро $0{,}5$). Значит, эта коробка строго не помещается.
\end{remark}

\section{Формальная проверка}\label{sec:lean}

Сопровождающий репозиторий \url{https://github.com/MikhailLevit1987/box-in-box} (архив в Zenodo: \href{https://doi.org/10.5281/zenodo.22974798}{doi:10.5281/zenodo.22974798}) содержит формализацию теоремы~\ref{thm:main} и утверждений
замечания~\ref{rem:strict} в Lean~4 (Lean v4.34.0-rc2, Mathlib), около 1800 строк в семи файлах:
\begin{itemize}
\item \texttt{Reduction.lean}, \texttt{Statement.lean}: определения, лемма~\ref{lem:widths} (через
теорему Мазура--Улама), лемма~\ref{lem:strip}(a), предложение~\ref{prop:suff};
\item \texttt{ZeroEntry.lean}: лемма~\ref{lem:strip}(b) (рассуждением с хордой вместо вогнутости),
предложение~\ref{prop:zero};
\item \texttt{Motzkin.lean}: теорема об альтернативах (через отделимость по Хану--Банаху);
\item \texttt{Descent.lean}: раздел~\ref{sec:descent};
\item \texttt{Necessity.lean}: раздел~\ref{sec:proof} и главная теорема \texttt{fits\_iff\_NC};
\item \texttt{Strict.lean}: раздел~\ref{sec:strict}, теоремы \texttt{fitsStrict\_iff\_shrink} и
\texttt{fitsStrict\_iff\_NCs} (с двумя дополнительными определениями: открытый параллелепипед и
строгое вложение).
\end{itemize}
Формулировка и два определения, от которых зависит её смысл, дословно:
\begin{verbatim}
def box (q : Fin 3 → ℝ) : Set E3 := {x | ∀ i, 0 ≤ x i ∧ x i ≤ q i}
def Fits (p q : Fin 3 → ℝ) : Prop := ∃ f : E3 ≃ᵢ E3, f '' box q ⊆ box p
theorem fits_iff_NC (p q : Fin 3 → ℝ) (hp : ∀ i, 0 < p i) (hq : ∀ i, 0 ≤ q i) :
    Fits p q ↔ NC p q
\end{verbatim}
Здесь \texttt{E3}~--- евклидово пространство $\R^3$, \texttt{box q}~--- множество точек $x$ с
$0\le x_i\le q_i$, а \texttt{Fits p q} означает, что некоторая изометрия $f$ пространства $\R^3$
переводит \texttt{box q} внутрь \texttt{box p}. Критерий \texttt{NC} дословно совпадает с
определением~\ref{def:NC} (с индексами $0,1,2$). Команда \texttt{\#print axioms Boxes.fits\_iff\_NC}
выдаёт только стандартные аксиомы \texttt{propext}, \texttt{Classical.choice}, \texttt{Quot.sound}.

\begin{thebibliography}{99}
\bibitem{AlexanderWetzel2005} R.~Alexander, J.~E.~Wetzel, Wafer in a box, \emph{Math. Mag.} 78 (2005) 214--220.
\bibitem{Carver1925} W.~B.~Carver, Solution of Problem 3036, \emph{Amer. Math. Monthly} 32 (1925) 47--49.
\bibitem{Carver1957} W.~B.~Carver, Solution of Problem E1225, \emph{Amer. Math. Monthly} 64 (1957) 114--116.
\bibitem{Garnett1923} F.~M.~Garnett, Problem 3036, \emph{Amer. Math. Monthly} 30 (1923) 337.
\bibitem{JerrardWetzel2004} R.~P.~Jerrard, J.~E.~Wetzel, Prince Rupert's rectangles, \emph{Amer. Math. Monthly} 111 (2004) 22--31.
\bibitem{JSSW2006} R.~Jerrard, J.~Schneider, R.~Smallberg, J.~Wetzel, Straw in a box, \emph{College Math. J.} 37 (2006) 93--102.
\bibitem{Kvant1999} Задача М1687 (Турнир городов, 1998), \emph{Квант}, 1999, №~6, с.~16.
\bibitem{MartinStephenson1989} R.~R.~Martin, P.~C.~Stephenson, Containment algorithms for objects in rectangular boxes, in: \emph{Theory and Practice of Geometric Modeling}, Springer, 1989, 307--325.
\bibitem{Motzkin1936} T.~S.~Motzkin, \emph{Beitr\"age zur Theorie der linearen Ungleichungen}, Dissertation, Basel, 1936.
\bibitem{Post1993} K.~A.~Post, Triangle in a triangle: on a problem of Steinhaus, \emph{Geom. Dedicata} 45 (1993) 201--225.
\bibitem{Schrijver1986} A.~Schrijver, \emph{Theory of Linear and Integer Programming}, Wiley, 1986.
\bibitem{vanSwinden1816} J.~H.~van Swinden, \emph{Grondbeginsels der Meetkunde}, 2nd ed., Amsterdam, 1816.
\bibitem{Wetzel2000} J.~E.~Wetzel, Rectangles in rectangles, \emph{Math. Mag.} 73 (2000) 204--211.
\end{thebibliography}

\end{document}
